% ===== PRÉAMBULE CANONIQUE — à coller VERBATIM en tête de CHAQUE .tex =====
\documentclass[11pt,a4paper]{article}
\usepackage[utf8]{inputenc}\usepackage[T1]{fontenc}\usepackage{lmodern}
\usepackage[french]{babel}
\usepackage{amsmath,amssymb,mathtools}\usepackage{icomma}
\usepackage{siunitx}\sisetup{locale=FR}  % \qty{}{}, \si{}, \ang{}
\usepackage{xcolor}\usepackage{colortbl}
\usepackage{tikz}\usepackage{pgfplots}\pgfplotsset{compat=1.18}
\usepackage[siunitx,europeanresistors]{circuitikz} % circuits (élec.)
\usepackage[most]{tcolorbox}\usepackage{enumitem}\usepackage{array,booktabs}
\usepackage[a4paper,margin=2.2cm]{geometry}
\usetikzlibrary{arrows.meta,calc,angles,quotes,positioning,patterns,%
  backgrounds,shapes.geometric,decorations.pathmorphing,decorations.markings,intersections}
\definecolor{primary}{RGB}{222,49,99}\definecolor{primarydark}{RGB}{150,22,55}
\definecolor{defcol}{RGB}{0,114,178}\definecolor{methcol}{RGB}{52,92,148}
\definecolor{excol}{RGB}{0,135,90}\definecolor{attncol}{RGB}{200,40,55}
\tcbset{boxbase/.style={enhanced,breakable,boxrule=0.9pt,arc=2.5pt,
  left=3mm,right=3mm,top=2.3mm,bottom=2.3mm,fonttitle=\bfseries,coltitle=white}}
% --- boîtes TUTORIEL (noms EXACTS, ne pas renommer) ---
\newtcolorbox{cle}[1][]{boxbase,colback=defcol!6,colframe=defcol,
  title={\textsc{À retenir}\ifx\relax#1\relax\relax\else\textmd{ : \itshape #1}\fi}}
\newtcolorbox{methode}[1][]{boxbase,colback=methcol!7,colframe=methcol,sharp corners,
  title={\textsc{Méthode}\ifx\relax#1\relax\relax\else\textmd{ --- \itshape #1}\fi}}
\newtcolorbox{exemplebox}[1][]{boxbase,colback=excol!5,colframe=excol,
  title={\textsc{Exemple}\ifx\relax#1\relax\relax\else\textmd{ : \itshape #1}\fi}}
\newtcolorbox{piege}[1][]{boxbase,colback=attncol!6,colframe=attncol,
  title={\textsc{Piège}\ifx\relax#1\relax\relax\else\textmd{ : \itshape #1}\fi}}
\newtcolorbox{remarque}[1][]{boxbase,colback=black!4,colframe=black!45,
  title={\textsc{Remarque}\ifx\relax#1\relax\relax\else\textmd{ : \itshape #1}\fi}}
% --- boîtes EXERCICE / CORRIGÉ (noms EXACTS, auto-comptées) ---
\newtcolorbox[auto counter]{exo}[1][]{boxbase,colback=primary!4,colframe=primary,
  title={\textsc{Exercice}~\thetcbcounter\ifx\relax#1\relax\relax\else\textmd{ --- \itshape #1}\fi}}
\newtcolorbox[auto counter]{corrige}[1][]{boxbase,colback=excol!4,colframe=excol,
  title={\textsc{Corrigé}~\thetcbcounter\ifx\relax#1\relax\relax\else\textmd{ --- \itshape #1}\fi}}
\newcommand{\vv}[1]{\vec{#1}}\newcommand{\ds}{\displaystyle}
\newcommand{\R}{\mathbb{R}}\newcommand{\N}{\mathbb{N}}\newcommand{\Z}{\mathbb{Z}}
% (un \usepackage{fancyhdr}+en-tête est possible ; il est ignoré côté web)
% =========================================================================
\begin{document}
\begin{exo}[à quelle distance retrouve-t-on le champ terrestre ?]
Une ligne électrique transporte un courant $I=\qty{100}{\ampere}$. À quelle distance $d$ du fil
l'intensité du champ magnétique est-elle égale à celle du champ terrestre
$B_T=\qty{5e-5}{\tesla}$ ? (On prend $\mu_0/2\pi=2\times10^{-7}$.)
\end{exo}

\begin{exo}[quel courant pour un champ donné ?]
On veut créer, à la distance $d=\qty{2}{\centi\metre}$ d'un fil rectiligne, un champ magnétique
d'intensité $B=\qty{1e-4}{\tesla}$. Quel courant $I$ faut-il faire passer dans le fil ?
\end{exo}

\begin{exo}[bobine plate de $N$ spires]
Une bobine plate de $N=50$ spires jointives, de rayon $r=\qty{5}{\centi\metre}$, est parcourue par
un courant $I=\qty{2}{\ampere}$.
\begin{enumerate}[label=\alph*)]
  \item Écris la formule du champ au centre d'une bobine de $N$ spires.
  \item Calcule l'intensité du champ au centre.
\end{enumerate}
\end{exo}

\begin{exo}[dimensionner un solénoïde]
On souhaite qu'un solénoïde de longueur $L=\qty{50}{\centi\metre}$ produise un champ interne
$B=\qty{1e-2}{\tesla}$ lorsqu'il est parcouru par $I=\qty{5}{\ampere}$. Combien de spires $N$
faut-il enrouler ? (On prend $\mu_0=\qty{1.26e-6}{}$.)
\end{exo}

\begin{exo}[deux bobines, même densité de spires]
On compare deux solénoïdes parcourus par le \textbf{même} courant $I$ :
\begin{center}
\renewcommand{\arraystretch}{1.2}
\begin{tabular}{@{}lcc@{}}
\toprule
 & Solénoïde 1 & Solénoïde 2 \\ \midrule
Nombre de spires $N$ & $100$ & $150$ \\
Longueur $L$ & \qty{4}{\centi\metre} & \qty{6}{\centi\metre} \\ \bottomrule
\end{tabular}
\end{center}
\begin{enumerate}[label=\alph*)]
  \item Calcule la densité linéique de spires $n=N/L$ pour chaque solénoïde.
  \item Lequel crée le champ interne le plus intense ? Justifie.
\end{enumerate}
\end{exo}

\end{document}
